% UFF 1997 — 15 questões MCQ — Química (excluídas 4 questões com imagem essencial)

---
tipo: Múltipla Escolha
dificuldade: Média
nivel: médio
disciplina: Química
assunto: Tabela Periódica e Propriedades Periódicas
gabarito: E
tags: [raio atômico, energia de ionização, afinidade eletrônica, propriedades periódicas]
fonte: concurso
concurso: UFF
banca: UFF
ano: 1997
numero: 1
---
\question
Analisando-se a classificação periódica dos elementos químicos, pode-se afirmar que:
\begin{choices}
\choice O raio atômico do nitrogênio é maior que o do fósforo.
\choice A afinidade eletrônica do cloro é menor que a do fósforo.
\choice O raio atômico do sódio é menor que o do magnésio.
\choice A energia de ionização do alumínio é maior que a do enxofre.
\CorrectChoice A energia de ionização do sódio é maior que a do potássio.
\end{choices}

---
tipo: Múltipla Escolha
dificuldade: Fácil
nivel: médio
disciplina: Química
assunto: Substâncias e Misturas
gabarito: D
tags: [estados físicos, soluções, misturas, classificação de soluções]
fonte: concurso
concurso: UFF
banca: UFF
ano: 1997
numero: 2
---
\question
São dadas as soluções:
\begin{enumerate}
\item argônio dissolvido em nitrogênio;
\item dióxido de carbono dissolvido em água;
\item etanol dissolvido em acetona;
\item mercúrio dissolvido em ouro.
\end{enumerate}
Estas soluções, à temperatura ambiente, são classificadas de acordo com seu estado físico em, respectivamente:
\begin{choices}
\choice líquida, líquida, gasosa, líquida
\choice gasosa, gasosa, líquida, sólida
\choice líquida, gasosa, líquida, líquida
\CorrectChoice gasosa, líquida, líquida, sólida
\choice líquida, gasosa, líquida, sólida
\end{choices}

---
tipo: Múltipla Escolha
dificuldade: Média
nivel: médio
disciplina: Química
assunto: Estequiometria
gabarito: A
tags: [reagente limitante, estequiometria, amônia, cloreto de amônio, óxido de cálcio]
fonte: concurso
concurso: UFF
banca: UFF
ano: 1997
numero: 3
---
\question
Amônia gasosa pode ser preparada pela seguinte reação balanceada:
$$\ce{CaO(s) + 2 NH4Cl(s) ->}$$
$$\ce{-> 2 NH3(g) + H2O(g) + CaCl2(s)}$$
Se 112,0 g de óxido de cálcio e 224,0 g de cloreto de amônio forem misturados, então a quantidade máxima, em gramas, de amônia produzida será, aproximadamente:

Dados: massas molares $\ce{CaO} = 56$ g/mol; $\ce{NH4Cl} = 53{,}5$ g/mol; $\ce{NH3} = 17$ g/mol.
\begin{choices}
\CorrectChoice 68,0
\choice 34,0
\choice 71,0
\choice 36,0
\choice 32,0
\end{choices}

---
tipo: Múltipla Escolha
dificuldade: Fácil
nivel: médio
disciplina: Química
assunto: Reações Químicas e Balanceamento
gabarito: D
tags: [tipos de reação, reação de substituição, síntese, decomposição]
fonte: concurso
concurso: UFF
banca: UFF
ano: 1997
numero: 4
---
\question
Considere as reações:
\begin{enumerate}
\item $\ce{4 Li + O2 -> 2 Li2O}$
\item $\ce{2 Li + MgO -> Li2O + Mg}$
\item $\ce{2 KClO3 -> 2 KCl + 3 O2}$
\item $\ce{Cl2 + 2 FeCl2 -> 2 FeCl3}$
\item $\ce{H2O2 -> H2O + 1/2 O2}$
\end{enumerate}
A reação de substituição é indicada por:
\begin{choices}
\choice I
\choice IV
\choice III
\CorrectChoice II
\choice V
\end{choices}

---
tipo: Múltipla Escolha
dificuldade: Média
nivel: médio
disciplina: Química
assunto: Pilhas e Células Galvânicas
gabarito: B
tags: [pilha galvânica, potencial padrão, eletrodo, eletroquímica]
fonte: concurso
concurso: UFF
banca: UFF
ano: 1997
numero: 5
---
\question
Em uma pilha galvânica, um eletrodo é cobre imerso em solução de $\ce{Cu^{2+}}$ 1,0 mol/L e o outro é prata imerso em solução de $\ce{Ag+}$ 1,0 mol/L.

Dados --- potenciais-padrão de redução a 25°C:
$$\ce{Cu^{2+} + 2 e^{-} -> Cu^0} \qquad E^\circ = 0{,}34 \text{ V}$$
$$\ce{Ag^{+} + e^{-} -> Ag^0} \qquad E^\circ = 0{,}80 \text{ V}$$
O potencial padrão da célula para esta pilha é:
\begin{choices}
\choice 1,14 V
\CorrectChoice 0,46 V
\choice 1,26 V
\choice 1,94 V
\choice 0,16 V
\end{choices}

---
tipo: Múltipla Escolha
dificuldade: Fácil
nivel: médio
disciplina: Química
assunto: Forças Intermoleculares
gabarito: C
tags: [forças intermoleculares, dipolo temporário, dipolo permanente, ligação de hidrogênio]
fonte: concurso
concurso: UFF
banca: UFF
ano: 1997
numero: 6
---
\question
Considere as seguintes interações:
\begin{enumerate}
\item $\ce{CH4} \cdots \ce{CH4}$
\item $\ce{HBr} \cdots \ce{HBr}$
\item $\ce{CH3OH} \cdots \ce{H2O}$
\end{enumerate}
As forças intermoleculares predominantes que atuam nas interações I, II e III são, respectivamente:
\begin{choices}
\choice ligação de hidrogênio, dipolo temporário, dipolo permanente
\choice ligação de hidrogênio, ligação de hidrogênio, dipolo temporário
\CorrectChoice dipolo temporário, dipolo permanente, ligação de hidrogênio
\choice dipolo temporário, ligação de hidrogênio, dipolo permanente
\choice dipolo permanente, ligação de hidrogênio, dipolo temporário
\end{choices}

---
tipo: Múltipla Escolha
dificuldade: Fácil
nivel: médio
disciplina: Química
assunto: Concentração Comum e Molar
gabarito: A
tags: [concentração molar, molaridade, vitamina C, ácido ascórbico]
fonte: concurso
concurso: UFF
banca: UFF
ano: 1997
numero: 7
---
\question
Dissolve-se 8,8 g de ácido ascórbico (vitamina C, $\ce{C6H8O6}$) em água suficiente para preparar 125 mL de solução.

A concentração molar deste componente na solução é:

Dado: massa molar $\ce{C6H8O6} = 176$ g/mol.
\begin{choices}
\CorrectChoice 0,40 mol/L
\choice 0,80 mol/L
\choice 0,10 mol/L
\choice 0,20 mol/L
\choice 1,00 mol/L
\end{choices}

---
tipo: Múltipla Escolha
dificuldade: Média
nivel: médio
disciplina: Química
assunto: Diluição e Mistura de Soluções
gabarito: C
tags: [diluição, concentração em massa, porcentagem em peso, solução]
fonte: concurso
concurso: UFF
banca: UFF
ano: 1997
numero: 8
---
\question
A solução de um certo sal tem a concentração de $30\%$ em peso. A massa de água necessária para diluí-la a $20\%$ em peso é:
\begin{choices}
\choice $25$ g
\choice $75$ g
\CorrectChoice $50$ g
\choice $100$ g
\choice $150$ g
\end{choices}

---
tipo: Múltipla Escolha
dificuldade: Fácil
nivel: médio
disciplina: Química
assunto: Equilíbrio de Solubilidade (Kps)
gabarito: D
tags: [Kps, produto de solubilidade, saturação, equilíbrio iônico]
fonte: concurso
concurso: UFF
banca: UFF
ano: 1997
numero: 9
---
\question
O seguinte equilíbrio ocorre em meio aquoso:
$$\ce{PbI2(s) <=> Pb^{2+}(aq) + 2 I^{-}(aq)} \qquad K_{ps}(\ce{PbI2}) = 8{,}3 \times 10^{-5}$$
Pode-se afirmar que:
\begin{choices}
\choice se $[\ce{Pb^{2+}}][\ce{I^{-}}]^2 = K_{ps}$, então a solução é insaturada.
\choice se $[\ce{Pb^{2+}}][\ce{I^{-}}]^2 > K_{ps}$, então a solução é saturada.
\choice se $[\ce{Pb^{2+}}][\ce{I^{-}}]^2 < K_{ps}$, então a solução é supersaturada.
\CorrectChoice se $[\ce{Pb^{2+}}][\ce{I^{-}}]^2 = K_{ps}$, então a solução é saturada.
\choice se $[\ce{Pb^{2+}}][\ce{I^{-}}]^2 > K_{ps}$, então a solução é insaturada.
\end{choices}

---
tipo: Múltipla Escolha
dificuldade: Fácil
nivel: médio
disciplina: Química
assunto: Estequiometria
gabarito: A
tags: [estequiometria, volume de gás, carbonato de cálcio, CNTP, decomposição]
fonte: concurso
concurso: UFF
banca: UFF
ano: 1997
numero: 10
---
\question
Para produzir 4,48 L de $\ce{CO2}$ nas CNTP, conforme a reação
$$\ce{CaCO3 -> CaO + CO2}$$
a quantidade necessária, em gramas, de $\ce{CaCO3}$ é:

Dado: massa molar $\ce{CaCO3} = 100$ g/mol.
\begin{choices}
\CorrectChoice 20,0
\choice 10,0
\choice 100,0
\choice 200,0
\choice 18,3
\end{choices}

---
tipo: Múltipla Escolha
dificuldade: Média
nivel: médio
disciplina: Química
assunto: Funções Orgânicas Oxigenadas
gabarito: B
tags: [fenol, funções orgânicas, radical arila, classificação orgânica]
fonte: concurso
concurso: UFF
banca: UFF
ano: 1997
numero: 11
---
\question
Para que a fórmula geral $\text{Y--OH}$ corresponda a uma função fenólica, $\text{Y}$ deve ser um radical:
\begin{choices}
\choice alcinila
\CorrectChoice arila
\choice cicloexila
\choice alcenila
\choice benzila
\end{choices}

---
tipo: Múltipla Escolha
dificuldade: Média
nivel: médio
disciplina: Química
assunto: Reatividade e Mecanismos de Reação
gabarito: E
tags: [substituição eletrofílica aromática, grupo amino, orientação orto-para, anel benzênico]
fonte: concurso
concurso: UFF
banca: UFF
ano: 1997
numero: 12
---
\question
O grupo amino ($\ce{-NH2}$), ligado ao anel benzênico, nas reações de substituição eletrofílica aromática é um orientador:
\begin{choices}
\choice apenas \textit{orto}
\choice \textit{meta} e \textit{para}
\choice apenas \textit{meta}
\choice \textit{orto} e \textit{meta}
\CorrectChoice \textit{orto} e \textit{para}
\end{choices}

---
tipo: Múltipla Escolha
dificuldade: Média
nivel: médio
disciplina: Química
assunto: Funções Orgânicas Oxigenadas
gabarito: A
tags: [acidez orgânica, fenol, álcool, éter, comparação de acidez]
fonte: concurso
concurso: UFF
banca: UFF
ano: 1997
numero: 13
---
\question
Considere os compostos:
\begin{enumerate}
\item Éter etílico
\item Fenol
\item $n$-Propanol
\end{enumerate}
Marque a opção que apresenta os compostos indicados em ordem \textbf{crescente} de acidez.
\begin{choices}
\CorrectChoice I, III, II
\choice II, I, III
\choice III, II, I
\choice I, II, III
\choice III, I, II
\end{choices}

---
tipo: Múltipla Escolha
dificuldade: Média
nivel: médio
disciplina: Química
assunto: Funções Orgânicas Oxigenadas
gabarito: B
tags: [carbonila, aldeído, cetona, hidratação de alcinos, funções orgânicas]
fonte: concurso
concurso: UFF
banca: UFF
ano: 1997
numero: 14
---
\question
Substâncias que contêm o grupo funcional carbonila podem ser obtidas a partir da seguinte reação:
\begin{choices}
\choice hidratação de alcenos
\CorrectChoice hidratação de alcinos
\choice hidrólise de éteres
\choice desidratação de álcoois
\choice hidrogenação de alcinos
\end{choices}

---
tipo: Múltipla Escolha
dificuldade: Fácil
nivel: médio
disciplina: Química
assunto: Reações Orgânicas
gabarito: D
tags: [esterificação, éster, ácido acético, isopropanol, reação orgânica]
fonte: concurso
concurso: UFF
banca: UFF
ano: 1997
numero: 15
---
\question
Na reação de $\ce{CH3COOH}$ com $\ce{CH3CH(OH)CH3}$, catalisada por ácido, além da água, resulta:
\begin{choices}
\choice isopropionato de etila
\choice acetato de propila
\choice 2-etil-1-propanol
\CorrectChoice acetato de isopropila
\choice 2-etil-2-propanol
\end{choices}
